\section{Introduction}
An Agent receives the instruction \emph{process the next support ticket}. It reads the ticket, checks an order and approval, consults a model, and submits a refund. Encryption hides message contents from a network observer, but the contacted services, traffic sizes and timing remain visible. The tool's answer then becomes part of a model request, which creates a second observable representation of the same task. Suppose a defender adds cover traffic. A second question follows: does that cover end at a proxy, execute a read, or change business state? Figure~\ref{fig:chain} follows information and actions through this example.

\begin{figure*}[t]
\centering
\begin{tikzpicture}[x=1cm,y=1cm,font=\small,>=Latex]
\tikzset{box/.style={draw,rounded corners,align=center,minimum height=.85cm,text width=2.4cm},
 note/.style={align=center,font=\footnotesize,text width=3.1cm}}
\node[box,fill=blue!5] (u) at (1.3,1.15) {User\\``Process the next ticket''};
\node[box,fill=blue!5] (a) at (5.1,1.15) {Agent\\plan, check, act, retry};
\node[box,fill=green!6] (m) at (9.2,2.6) {Model service\\next action / streamed output};
\node[box,fill=green!6] (t) at (9.2,-.35) {Tool service\\ticket, order, refund};
\node[box,fill=orange!8] (d) at (14.1,1.15) {Specialist Agent\\optional delegation};
\draw[<->,thick] (u)--node[above,yshift=15pt]{client link C}(a);
\draw[<->,thick] (a.north east)--node[above,sloped]{model link M}(m.west);
\draw[<->,thick] (a.south east)--node[below,sloped]{tool link T}(t.west);
\draw[<->,thick] (a.east)--node[above]{coordination A}(d.west);
\node[note] at (2.0,-.35) {\textbf{Transit observer}\\lengths, gaps, connections;\\destinations when available};
\node[note] at (13.7,-.35) {\textbf{Recipient / active service}\\received arguments;\\optionally selected responses};
\node[draw=black!45,fill=black!3,align=left,text width=15.7cm,inner sep=6pt] at (7.9,-1.65) {
\textbf{Two protection questions:} Where does tool feedback appear next? What work does privacy cover execute before it is discarded? The original task and each observable service remain part of the protection decision.};
\end{tikzpicture}
\caption{Running support-ticket example and observation rights. A tool result returns to the Agent and can enter the next model request. Each encrypted link exposes its own traffic shape; a recipient additionally sees the request addressed to it. Concurrent tickets may share a connection on either link. The arrows represent application communication, not a one-to-one mapping from tool events to packets.}
\label{fig:chain}
\end{figure*}


\paragraph{Why systematize now}
Agent deployments combine remote tool protocols, including the Model Context Protocol (MCP), with streamed model APIs~\cite{mcp2025transports,rfc9113}. In under three years, studies have recovered research queries from contacted domains~\cite{jeong2026wra}, recognized tasks from model traffic~\cite{li2026ant,pouryousef2026traffic}, fingerprinted hosted Agents~\cite{agentprint2025} and attributed models from tool-call behavior~\cite{zhao2026agenttrace,ediga2026trace}. The first defenses are research decoys, relays and shaping~\cite{jeong2026wra,dangol2026workflow}. Most of this work is preprint-stage, and no shared evaluation standard exists. Website fingerprinting (WF) went through the same phase, and many early attack and defense results were later reassessed under realistic populations and adapted attackers~\cite{dyer2012peekaboo,juarez2014critical,mathews2023critical}. An SoK at this point can set evaluation expectations before results pile up that cannot be compared. Its corpus is necessarily small, so we make every conclusion's source dependence explicit.

\paragraph{The question}
Which traffic-analysis evaluation conditions carry over to Agent execution, and which need additional execution conditions? We organize the answer around two boundaries: the services through which task information travels, and the application work that cover performs. Classical evaluations already address multiple observations and adaptive attacks; a design with a single protected link or a tunnel terminator supplies a narrower protection scope than an Agent executing work across services. We turn these distinctions into an evaluation procedure, then apply it to cover generated by a controller or by a second planner.

\textbf{Finding 1: protection scope follows the task across service links.} WRA recovers research queries from tool destinations, while ANT and a cloud benchmark recognize task labels on model connections~\cite{jeong2026wra,li2026ant,pouryousef2026traffic}. A contacted tool additionally reads its arguments, and PrivacyTrace shows that these arguments raise attribute inference well above a schema prior~\cite{duan2026privacytrace}. Tool results are forwarded into later model requests, so the same task is observable on links that different parties control, and a tool-side defense does not cover the model link. A defense claim must therefore name every view its attacker retains (Section~\ref{sec:authorization}). Our forwarding pilot only calibrates this consequence; its primary tests do not reject (Section~\ref{sec:live-study}).

\textbf{Finding 2: executed cover makes utility part of the privacy result.} Classical cover terminates at a tunnel endpoint~\cite{apthorpe2019stp}. An Agent decoy can execute queries or update state. The published research-decoy result lowers query recovery from .770 to .565, but its utility scores belong to a separate destination-visibility intervention~\cite{jeong2026wra}. Across the four mapped query/workflow privacy settings, adapted attacks appear in two and matched task utility in one, and no setting reports both (Section~\ref{sec:reporting}). These mapped records provide no setting with both adapted inference and matched task utility, so their joint evidence remains incomplete.

\textbf{Finding 3: the cover's dispatcher is part of its observable behavior.} In a new 3,000-episode study, immediate controller cover can be stripped without cover labels: task recognition returns to 112/120 and 108/120 for two fixed-weight planners. Live historical-gap dispatch blocks that simple shortcut but leaves context-removal and joint-decoding attacks effective; under concurrent dispatch their gains are 27.5 and 31.7 points after Holm correction. With a parallel shadow planner, neither validation selector chooses a timing attack; model tokens approximately double, while a size/destination RF still recognizes 70 and 64 tasks. This is a measured dispatch tradeoff, with task and state outcomes on the same runs (Section~\ref{sec:cover-study}).

\paragraph{Contributions}
\begin{itemize}
\item \textbf{An access-aware evidence map and transfer matrix}: 44 settings from 15 Agent studies, connected to 22 traffic-analysis anchors. Each result retains its target, observation rights, population and supplied boundaries (Sections~\ref{sec:scope}--\ref{sec:correspondence}).
\item \textbf{Two execution-aware syntheses}, of feedback exposure and cover termination, yielding six defense layers, a four-step evaluation procedure and source-withdrawal results (Sections~\ref{sec:evidence}--\ref{sec:defenses}, \ref{sec:reporting}--\ref{sec:agenda}).
\item \textbf{A worked technical case on actual dispatch}: 3,000 fixed-weight planner episodes compare immediate, history-based and planner-generated cover, with adapted inference, goal/state outcomes and resource costs. Earlier receipt analyses and the separate forwarding pilot retain their own results and denominators (Section~\ref{sec:experiment}).
\end{itemize}
